\documentclass[aps,pre,superscriptaddress,floatfix,longbibliography]{revtex4-2}

\usepackage{amsmath}
\usepackage{amssymb}
\usepackage{bm}
\usepackage{graphicx}

\begin{document}

\title{Supplementary Material}

\maketitle

\section{Monod-Wyman-Changeux model for heteromeric receptors}
\label{sec:mwc}

The single-channel opening probability used in the main text extends the
Monod-Wyman-Changeux (MWC) treatment of ligand-gated ion channels of
Ref.~\cite{einav_monod-wyman-changeux_2017}, which derives the opening curve
of a homomeric channel.%, all $k_\text{sub}$ subunits identical. 
Our
receptors are heteromers: a ring of $k_\text{sub}$ subunits
$r=(u_1,\dots,u_{k_\text{sub}})$, each drawn from a pool of $n_\text{genes}$
possible gene products and chemically distinct in general. The ligand does not bind an isolated
subunit but the pocket formed at the interface between two
neighbouring subunits (Fig.~\ref{fig:pentamer}). Each subunit is structurally
asymmetric, presenting two different faces to its two neighbours. We label
them "$+$" and "$-$". The ring therefore has exactly $k_\text{sub}$ pockets,
pocket $i$ formed by the "$+$" face of $u_i$ meeting the "$-$" face of
$u_{i+1}$. Because the two faces of a subunit
are chemically distinct, the pocket formed by $(u_i^+,u_j^-)$ is generally
not the same as the one formed by $(u_j^+,u_i^-)$, see Fig.~\ref{fig:pentamer}. For a given pair of
subunit species there are two possible pocket types, depending on which one
contributes its "$+$" face. Because the reference does not itself treat
heteromers or interfacial pockets, we derive the model below.

\begin{figure}[t]
\centering
\includegraphics[width=0.6\linewidth]{figure/pentamer_formula.pdf}
\caption{Heteropentameric receptor ring, subunits shown with their two
structurally distinct faces. Each of the ligand pockets sits
at the interface between the "$+$" face of one subunit and the "$-$" face of its neighbour.
As a result, homopentamer have only one type of pocket (left-hand side), whereas heteropentamer have two types of pockets (right-hand side). }
\label{fig:pentamer}
\end{figure}

\paragraph{Binding energy and dissociation constant.} For pocket $i$ in
conformation $\alpha\in\{o,c\}$, ligand binding is the reaction
$P_i^\alpha + L \rightleftharpoons P_i^\alpha L$, with dissociation constant
$K_\alpha^{(i,\ell)} = [P_i^\alpha][L]/[P_i^\alpha L]$. Each species'
equilibrium concentration is set by its Boltzmann weight,
$[s]\propto e^{-\beta E_s}$, so
\begin{equation}
K_\alpha^{(i,\ell)} = \frac{e^{-\beta E(P_i^\alpha)}\,e^{-\beta E(L)}}{e^{-\beta E(P_i^\alpha L)}}
= e^{\beta \Delta E_\alpha^{(i,\ell)}},
\qquad
\Delta E_\alpha^{(i,\ell)} \equiv E(P_i^\alpha L) - E(P_i^\alpha) - E(L),
\label{eq:K_from_E}
\end{equation}
i.e. the dissociation constant is set by the binding free energy
$\Delta E_\alpha^{(i,\ell)}$ of the ligand to that specific pocket (we set
$\beta=1$, energies in units of $k_BT$, throughout). It is this energy,
rather than $K_\alpha^{(i,\ell)}$ itself, that we later parametrize directly
in the latent space (Sec.~\ref{sec:kernel}).

\paragraph{Derivation.} The channel is a two-state allosteric system: it is
found either fully open ($C_o$) or fully closed ($C_c$), with intrinsic
(ligand-free) energies $E_o$ and $E_c$. We set $E_o=0$ and write
$\epsilon_r = E_c-E_o$ for the leakiness. Within either conformation
$\alpha$, ligand binds each of the $k_\text{sub}$ pockets independently, so
pocket $i$ is either empty (statistical weight $1$) or ligand-bound (weight
$c/K_\alpha^{(i,\ell)}$). Summing over the $2^{k_\text{sub}}$
ligand-occupancy microstates of conformation $\alpha$ then factorizes into a
product over the independent pockets, the conformation's binding polynomial:
\begin{equation}
Z_\alpha(c,\ell) = \prod_{i=1}^{k_\text{sub}}\left(1+\frac{c}{K_\alpha^{(i,\ell)}}\right).
\label{eq:binding_poly}
\end{equation}
For a homomer, all pockets are chemically identical and
Eq.~\eqref{eq:binding_poly} collapses to the usual $(1+c/K_\alpha)^{k_\text{sub}}$
form of Ref.~\cite{einav_monod-wyman-changeux_2017}. For a heteromer the
pockets differ and the product must be kept in full. This is where the sum
of individual pocket interaction energies enters the model. Weighting each
conformation's branch by its Boltzmann factor, the opening probability is the
fraction of the total statistical weight found open,
\begin{equation}
p_o^{(r,\ell)}(c) = \frac{Z_o(c,\ell)}{Z_o(c,\ell) + e^{-\epsilon_r} Z_c(c,\ell)}
= \frac{\prod_{i=1}^{k_\text{sub}} (1+c/K_o^{(i,\ell)})}{\prod_{i=1}^{k_\text{sub}} (1+c/K_o^{(i,\ell)}) + e^{-\epsilon_r}\prod_{i=1}^{k_\text{sub}}(1+c/K_c^{(i,\ell)})},
\label{eq:po_heteromer}
\end{equation}
which is the heteromeric extension used in the main text.

\paragraph{Rescaled affinity $\tilde K_o$.} Fitting Eq.~\eqref{eq:po_heteromer}
directly is degenerate: Ref.~\cite{einav_monod-wyman-changeux_2017} shows, via
the Hessian of $p_o$ with respect to $(\epsilon_r,\Delta E_o,\Delta E_c)$,
that one eigendirection is far softer ("sloppy") than the other two, so that
dose-response data alone cannot separate $\epsilon_r$ from $K_o$. Only the
combination $e^{-\epsilon_r/2}K_o$ is identifiable. We therefore absorb the
leakiness into a rescaled open-state affinity rather than carry $\epsilon_r$
as an independent parameter. Distributing $\epsilon_r=\sum_{i=1}^{k_\text{sub}}\epsilon_i$
evenly among the $k_\text{sub}$ pockets in log space,
\begin{equation}
\tilde K_o^{(i,\ell)} = K_o^{(i,\ell)}\, e^{-\epsilon_r/k_\text{sub}},
\qquad\text{so that}\qquad
e^{-\epsilon_r}\prod_{i=1}^{k_\text{sub}} K_o^{(i,\ell)} = \prod_{i=1}^{k_\text{sub}} \tilde K_o^{(i,\ell)},
\end{equation}
and keeping only the identifiable, non-sloppy combination reduces
Eq.~\eqref{eq:po_heteromer} to
\begin{equation}
p_o^{(r,\ell)}(c) = \frac{\prod_{i=1}^{k_\text{sub}} (c/\tilde K_o^{(i,\ell)})}{\prod_{i=1}^{k_\text{sub}}(c/\tilde K_o^{(i,\ell)}) + \prod_{i=1}^{k_\text{sub}}(1+c/K_c^{(i,\ell)})}.
\label{eq:po_rescaled}
\end{equation}
$\tilde K_o^{(i,\ell)}$, not $K_o^{(i,\ell)}$ and $\epsilon_r$ separately, is
therefore the physically meaningful open-state affinity of pocket $i$, and is
what we use throughout.

\paragraph{Threshold reduction.} We further reduce
Eq.~\eqref{eq:po_rescaled} to a strict activation threshold, the
half-activation concentration $c^*$ at which $p_o=1/2$. This is a choice
of parsimony rather than a claim about receptor physics: it removes the
closed-state affinities $K_c^{(i,\ell)}$ from the model entirely, leaving a
single free affinity per pocket-ligand pair, the minimal parametrization
needed to give each pocket a tunable, combinable threshold.

To see why $K_c^{(i,\ell)}$ drops out, set $p_o(c^*)=1/2$ in
Eq.~\eqref{eq:po_rescaled}, which requires the two branches to balance,
\begin{equation}
\prod_{i=1}^{k_\text{sub}} \frac{c^*}{\tilde K_o^{(i,\ell)}} = \prod_{i=1}^{k_\text{sub}} \left(1+\frac{c^*}{K_c^{(i,\ell)}}\right).
\label{eq:balance}
\end{equation}
A functional ion channel is "good" in the MWC sense: the closed conformation
barely binds the ligand, $K_c^{(i,\ell)} \gg K_o^{(i,\ell)}$ for every
pocket. Since the threshold itself sits at the scale of the open-state
affinities, $c^*\sim\tilde K_o^{(i,\ell)} \ll K_c^{(i,\ell)}$, each factor on
the right-hand side of Eq.~\eqref{eq:balance} is close to its ligand-free
value, $1+c^*/K_c^{(i,\ell)} \approx 1$, and the right-hand side reduces to
$1$. Eq.~\eqref{eq:balance} then gives $(c^*)^{k_\text{sub}} =
\prod_{i=1}^{k_\text{sub}}\tilde K_o^{(i,\ell)}$, i.e. the threshold
$c^*_{r,\ell}$ is the geometric mean of the pockets' rescaled affinities,
\begin{equation}
c^*_{r,\ell} = \Bigl(\prod_{i=1}^{k_\text{sub}} \tilde K_o^{(i,\ell)}\Bigr)^{1/k_\text{sub}}
\;\;\Longleftrightarrow\;\;
\ln c^*_{r,\ell} = \frac{1}{k_\text{sub}}\sum_{i=1}^{k_\text{sub}} E_o^{(i,\ell)},
\qquad E_o^{(i,\ell)} \equiv \ln\tilde K_o^{(i,\ell)},
\label{eq:ec50}
\end{equation}
with the receptor activating once $c$ exceeds this threshold.

\section{Saturating affinity kernel}
\label{sec:kernel}

Section~\ref{sec:mwc} leaves $E_o^{(i,\ell)}$, the open-state binding energy
of pocket $i$, unspecified. We now give it a concrete form in the
morphochemical latent space introduced in Fig.~1\textbf{A} of the main text,
in a way that respects the two-face asymmetry of Sec.~\ref{sec:mwc}.

\paragraph{Two chemistries per subunit.} A ligand $\ell$ carries a fixed
latent coordinate $\mathbf v_\ell\in\mathbb R^D$. Because a subunit's "$+$"
and "$-$" faces make chemically distinct contacts (Fig.~\ref{fig:pentamer}),
a single latent vector per subunit cannot represent both: subunit $u$ is
instead assigned two independent latent coordinates, $\mathbf v_u^+$ and
$\mathbf v_u^-\in\mathbb R^D$, one per face-chemistry, each with its own pair
of learnable energy parameters $(E_\text{base}^{u,+},E_\text{max}^{u,+})$ and
$(E_\text{base}^{u,-},E_\text{max}^{u,-})$ — the meaning of these parameters
is given below.

\paragraph{Pocket embedding.} Pocket $i$, formed by the "$+$" face of $u_i$
meeting the "$-$" face of $u_{i+1}$ (Sec.~\ref{sec:mwc}), inherits a
chemistry that is a combination of the two contributing faces. We represent
it by the midpoint of their latent coordinates and of their energy
parameters,
\begin{equation}
\mathbf v_i = \tfrac{1}{2}\bigl(\mathbf v_{u_i}^{+} + \mathbf v_{u_{i+1}}^{-}\bigr),
\qquad
E_\text{base}^{(i)} = \tfrac{1}{2}\bigl(E_\text{base}^{u_i,+} + E_\text{base}^{u_{i+1},-}\bigr),
\qquad
E_\text{max}^{(i)} = \tfrac{1}{2}\bigl(E_\text{max}^{u_i,+} + E_\text{max}^{u_{i+1},-}\bigr).
\label{eq:pocket_embed}
\end{equation}
Averaging keeps the pocket's descriptors in the same $D$-dimensional space as
a single face, so the ligand kernel below needs only one distance per
pocket; the average is taken \emph{before} the kernel is applied (Eq.~\eqref{eq:Eo_kernel} below), not after, since the kernel is nonlinear in
$\mathbf v_i$.

\paragraph{Kernel.} The open-state binding energy of pocket $i$ for ligand
$\ell$ is a Gaussian radial basis function (RBF) of the latent distance
$d_i = \|\mathbf v_i - \mathbf v_\ell\|$,
\begin{equation}
E_o^{(i,\ell)} = E_\text{base}^{(i)} + E_\text{max}^{(i)}\left(1-\exp\!\left(-\frac{d_i^2}{\lambda^2}\right)\right),
\label{eq:Eo_kernel}
\end{equation}
with $\lambda$ a single length scale shared by every pocket (a fixed
hyperparameter, not learned), while $E_\text{base}^{(i)}$ and
$E_\text{max}^{(i)}$ inherit the per-face learnable parameters through
Eq.~\eqref{eq:pocket_embed}. In terms of the threshold $c^*$ of
Sec.~\ref{sec:mwc}, $\ln c^*_\text{min}\equiv E_\text{base}^{(i)}$ is the
threshold at perfect match ($d_i=0$) and $\ln c^*_\text{max}\equiv
E_\text{base}^{(i)}+E_\text{max}^{(i)}$ is the threshold reached as
$d_i\to\infty$: $c^*$ therefore rises smoothly between these two bounds as
the ligand moves away from the pocket's optimum, saturating rather than
diverging, since a fully mismatched ligand should incur a bounded, not
unbounded, energetic penalty; for $d_i\ll\lambda$, Eq.~\eqref{eq:Eo_kernel}
reduces to the harmonic form $E_o^{(i,\ell)}\approx E_\text{base}^{(i)} +
(E_\text{max}^{(i)}/\lambda^2)\,d_i^2$.

Combined with Eq.~\eqref{eq:ec50}, the heteromer's activation threshold is
\begin{equation}
\ln c^*_{r,\ell} = \frac{1}{k_\text{sub}}\sum_{i=1}^{k_\text{sub}}
\left[E_\text{base}^{(i)} + E_\text{max}^{(i)}\left(1-\exp\!\left(-\frac{d_i^2}{\lambda^2}\right)\right)\right],
\end{equation}
so that a heteromer's threshold for a given ligand is set jointly by how
well that ligand matches each of its $k_\text{sub}$ interfacial pockets.
Fig.~\ref{fig:env}\textbf{B} shows the shape of $c^*$ as a function of the
latent distance $d$ that this kernel produces at the single-pocket level,
saturating between $c^*_\text{min}$ and $c^*_\text{max}$ as derived above.

\section{Generation of the environment and the simulation loop}
\label{sec:env}

\begin{figure}
    \includegraphics[width=1.\linewidth]{../paragraph/fig2/figure_composite.pdf}
    \caption{\label{fig:env}Pipeline for estimating the mutual information of the array.
    \textbf{(A)} Two-dimensional UMAP projection of the $D$-dimensional
    morphochemical latent space. Ligands (dots) are colored by chemical family,
    receptors (open circles) are numbered. The affinity of a receptor for a
    ligand is set by their distance $d = \|\mathbf{v}_u - \mathbf{v}_\ell\|$
    (red arrow).
    \textbf{(B)} Activation threshold $c^*$ as a function of $d$, saturating
    between $c^*_\text{min}$ at complete match and $c^*_\text{max}$ at complete
    mismatch.
    \textbf{(C)} Monte Carlo loop. The environment draws the per-receptor
    threshold vector $\mathbf{c}^*_{\mathbf{r},\ell}$ over the $L$ ligands and a
    mixture $\mathbf{c} = (n_1 c_1,\dots,n_L c_L)$, with presence
    $n_\ell \sim \text{Bernoulli}(p)$ and concentration
    $c_\ell \sim \text{LogNormal}(\mu,\sigma)$. The physics step turns this into
    a binary activation pattern $A = (1,0,\dots,1)$. Repeating over sensing
    events builds the empirical distribution $p(A)$, from which the mutual
    information $I(A;(\mathbf{c},\{\ell\})) = H(A)$ is computed. The receptor
    vectors $\{\mathbf{v}_u\}$ are then optimized to maximize it.}
\end{figure}

\paragraph{Fixed pool.} $n_\text{families}$ family prototypes are placed in
latent space with a target average inter-family distance $d_\text{fam}$.
$L$ ligands are then drawn once around their family's prototype,
$\mathbf v_\ell \sim \mathcal N(\mathbf v_f,\sigma_\text{shape})$, together
with per-ligand log-concentration parameters $(\mu_{c,\ell},\sigma_{c,\ell})$.
Fig.~\ref{fig:env}\textbf{A} shows this latent space, with ligands colored by
family and receptors placed among them. This fixed pool, together with the
kernel length scale $\lambda$ of Sec.~\ref{sec:kernel}, is the full set of
parameters defining the environment, summarized in Table~\ref{tab:env_params}.

\begin{table}[t]
\centering
\begin{tabular}{ll}
Symbol & Meaning \\
\hline
$L$ & number of ligands in the fixed pool \\
$n_\text{families}$ & number of chemical families \\
$D$ & latent space dimension \\
$\sigma_\text{shape}$ & intra-family spread \\
$d_\text{fam}$ & average inter-family distance \\
$\mu_S$ & mean number of ligands per sniff \\
$(\mu_{c,\ell},\sigma_{c,\ell})$ & log-concentration mean and spread of ligand $\ell$ \\
$\sigma_\text{noise}$ & observation noise \\
$\lambda$ & kernel length scale (Sec.~\ref{sec:kernel}) \\
\end{tabular}
\caption{Parameters defining the environment. Used in Sec.~\ref{sec:perfect}
to specify the regime in which the results are reported.}
\label{tab:env_params}
\end{table}

\paragraph{Sensing event.} At each sensing event, the number of ligands
present, $S$, is drawn from a zero-truncated Poisson distribution with rate
$\mu_S$, restricted to the pool size,
\begin{equation}
P(S=s) = \frac{\mu_S^s/s!}{\sum_{k=1}^{L}\mu_S^k/k!}, \qquad s=1,\dots,L,
\label{eq:ztp}
\end{equation}
so that every sniff contains at least one ligand and the mean sniff size is
set by $\mu_S$. Given $S$, the $S$ present ligands are chosen uniformly at
random, without replacement, from the pool of $L$, which we encode as a
presence mask $M=(M_1,\dots,M_L)\in\{0,1\}^L$, $M_\ell=1$ if ligand $\ell$ is
among the $S$ chosen, $0$ otherwise, so $\sum_\ell M_\ell = S$. Each present
ligand draws its own concentration, $c_\ell\sim\text{LogNormal}(\mu_{c,\ell},\sigma_{c,\ell})$,
and its own observation noise, $\mathbf v_{\text{obs},\ell}\sim\mathcal
N(\mathbf v_\ell,\sigma_\text{noise})$ with $\sigma_\text{noise}\ll\sigma_\text{shape}$,
representing docking or thermal fluctuations around its fixed pool
coordinate.

This mask carries no structure linking which particular ligands co-occur,
beyond the constraint that exactly $S$ of them are present. Other models of
the olfactory environment make the opposite choice and build in such
structure: through joint-occurrence correlations in
Refs.~\cite{zwicker_receptor_2016,tesileanu_adaptation_2019}, or through
explicit shared odor sources in Ref.~\cite{del_castillo_convergent_2026}. We
omit this here to keep the number of free environment parameters small.

\paragraph{Fixed versus optimized.} The environment of
Table~\ref{tab:env_params} is sampled once and held fixed throughout
training. The only quantities updated by gradient descent are the receptor
chemistry of Sec.~\ref{sec:kernel}: the per-face latent coordinates
$\mathbf v_u^{+},\mathbf v_u^{-}$ and energy parameters
$E_\text{base}^{u,\pm},E_\text{max}^{u,\pm}$ of every unit $u$, updated to
maximize the array's mutual information (main text). Fig.~\ref{fig:env}\textbf{C}
summarizes the resulting pipeline: the environment and the physics of
Secs.~\ref{sec:mwc} and \ref{sec:kernel} generate the activation pattern $A$ for
each sensing event, repeated sensing events build the empirical distribution
$p(A)$ from which $H(A)$ is estimated, and the receptor coordinates
$\{\mathbf v_u^{\pm}\}$ are updated to maximize it.

\paragraph{From single-ligand threshold to mixture drive.}
Sec.~\ref{sec:mwc} defines activation for a single ligand as a hard
threshold, $A_r=\Theta(\ln c-\ln c^*_{r,\ell})$. In a mixture, ligands
compete for the same receptor, and each contributes to activation in
proportion to how far its concentration sits above its own threshold,
$c_\ell/c^*_{r,\ell}$. Summing these occupancy ratios over the present
ligands defines the receptor's drive,
\begin{equation}
\text{drive}_r(\{c_\ell\}) = \sum_{\ell:\,M_\ell=1} \frac{c_\ell}{c^*_{r,\ell}}.
\label{eq:drive}
\end{equation}
This generalizes the single-ligand rule: half-activation is reached exactly
when $\text{drive}_r=1$, which reduces to $c=c^*_{r,\ell}$ for a single ligand.

\paragraph{Soft binning.} The hard threshold $\Theta$ has zero gradient
almost everywhere and cannot be optimized by backpropagation. We replace it
by a temperature-scaled sigmoid of the log-drive,
\begin{equation}
p(r\mid\{c_\ell\}) = \sigma\!\left(\frac{\ln\,\text{drive}_r(\{c_\ell\})}{T}\right).
\label{eq:soft_bin}
\end{equation}
Evaluating $\ln\,\text{drive}_r$ directly from Eq.~\eqref{eq:drive} is
numerically unstable: the ratios $c_\ell/c^*_{r,\ell}$ can span many
orders of magnitude, since concentrations are log-normally distributed and
$c^*_{r,\ell}$ ranges between $c^*_\text{min}$ and $c^*_\text{max}$, so summing
them directly in floating point can overflow or silently underflow the
smallest terms to zero. We instead compute the equivalent logsumexp form,
\begin{equation}
\ln\,\text{drive}_r(\{c_\ell\}) = \ln\sum_{\ell:\,M_\ell=1} \exp\bigl(\ln c_\ell - \ln c^*_{r,\ell}\bigr),
\label{eq:logsumexp}
\end{equation}
which factors out the largest term before exponentiating the rest, so every
exponentiated term stays of order $1$ or smaller and the sum is computed at
the working numerical precision. $T$ is annealed during training, from an
initial value down to a small target value. In the limit $T\to0$,
Eq.~\eqref{eq:soft_bin} goes to the discrete binary activation model defined
in the main text, $p(r\mid\{c_\ell\})\to A_r\in\{0,1\}$.

\section{The perfect-array regime}
\label{sec:perfect}

$H(A)$ can depend on every environment parameter of Table~\ref{tab:env_params}
and on $R$ itself. A direct comparison between homomer and heteromer arrays
is only meaningful in a regime where it does not. We look for a window of
parameter values in which $H(A)$ stops changing with their precise value.
Inside this window the homomer array reaches its architectural ceiling,
$H(A)=R$. We call such an array perfect.

In this same window, with the same environment and the same fitting
procedure, heteromer arrays fall short of $R$ (main text). This is the
point of the window. Since the homomer array already reaches $R$ throughout
it, neither the environment nor the optimization can explain a heteromer
shortfall. Any gap measured between heteromer and homomer arrays of the
same $R$ is therefore a property of the heteromeric array itself, its
subunit-sharing correlations (Sec.~\ref{sec:mwc}), not an artifact of the
environment or of how it is fit.

Even for the homomer array, $H(A)$ can fail to reach $R$ outside this
window, for three distinct reasons.

\begin{itemize}
\item \textbf{Statistical.} $A$ is a deterministic, noiseless function of
the mixture (Sec.~\ref{sec:env}), so $H(A)\le H(\text{mixture})$: a
deterministic map cannot create entropy. If the mixture carries little
entropy, too sparse or too few ligands, $H(\text{mixture})$ binds before
$R$ does.
\item \textbf{Geometric.} Reaching $H(A)=R$ needs $R$ receptors that
realize $R$ mutually decorrelated, half-active dichotomies of the ligand
pool. Whether they can is a question about the expressive capacity of the
receptor family in $D$ dimensions, not about $H(\text{mixture})$.
\item \textbf{Numerical.} The first two mechanisms are properties of the
physical model. The gradient-based procedure of Sec.~\ref{sec:env} can
separately fail to reach or measure $H(A)=R$, usually for reasons of
computational cost.
\end{itemize}

Table~\ref{tab:param_windows} lists, for each parameter, the range over
which $H(A)=R$ holds, and which mechanism bounds it.

\begin{table}[t]
\centering
\begin{tabular}{p{2.2cm}p{2.6cm}p{5.2cm}p{2.0cm}}
Parameter & Range & Justification & Mechanism \\
\hline
$\mu_S,\,L$ & not too small & Too sparse or too few ligands caps
$H(\text{mixture})$ below $R$. & Statistical \\[6pt]
$\rho=\dfrac{\sigma_\text{shape}\sqrt{D}}{\lambda}$ & $\sim[0.2,1]$ &
Below: ligands within a family share one receptor's ball. Above: each ball
covers at most one ligand. & Geometric \\[6pt]
$d_\text{fam}/\lambda$ & $\sim[0.5,1.5]$ & Below: neighbouring families'
balls overlap. Above: no ball bridges families. & Geometric \\[6pt]
$D$ & no upper bound & Larger $D$ never hurts capacity, only costs memory
and parameters. We use the smallest $D$ satisfying the $\rho$,
$d_\text{fam}/\lambda$ windows above. & Numerical \\
\end{tabular}
\caption{Parameter windows for the perfect-array regime
(Sec.~\ref{sec:perfect}).}
\label{tab:param_windows}
\end{table}

The geometric mechanism is a capacity question, classically studied for
linear classifiers, hyperplanes~\cite{cover_geometrical_1965}. The VC
dimension of a hyperplane in $\mathbb{R}^D$ is $D+1$: any $D+1$ points in
general position can be split into any of their $2^{D+1}$ dichotomies, but
this guarantee is lost beyond $D+1$ points.

Our receptors are not hyperplanes. Through the Gaussian kernel of
Sec.~\ref{sec:kernel}, each is a ``ball'' classifier: it responds to
ligands within roughly $\lambda$ of its position, not to one side of a
plane. Kernel classifiers of this kind are generally less constrained by
the ambient dimension than a hyperplane is. A nonlinear kernel implicitly
maps the ligands into a much higher, for a Gaussian kernel effectively
infinite, dimensional feature space, where separating them is easier. This
is the same reasoning Cover used to motivate feature expansion in the
first place.

The hyperplane bound is therefore a conservative reference point, not a
tight limit for our receptors. What we keep from the comparison is the
direction of the effect: increasing $D$ can only make the classifier at
least as capable, never less. This is why the smallest $D$ consistent with
the $\rho$ and $d_\text{fam}/\lambda$ windows costs nothing in capacity.

\section{Bracketing the joint Shannon entropy}
\label{sec:entropy}

Exact evaluation of $H(A) = -\sum_A p(A)\log_2 p(A)$ requires enumerating
all $2^R$ activation patterns. Estimating $p(A)$ reliably from $B$ samples
needs roughly one sample per pattern, $B\sim2^R$, and the intermediate
tensor built from these samples has $(2^R)^2=2^{2R}$ entries. At float32
precision this is $2^{2R+2}$ bytes. At $R=15$ this is $2^{32}$ bytes,
$4$ GiB, comfortable on an 80 GB A100. Each additional receptor quadruples
it. By $R=18$, three receptors later, it is $2^{38}$ bytes, $256$ GiB,
beyond any single GPU. We therefore need an estimator that does not
enumerate $2^R$ patterns once $R$ grows past the high teens.

The object we want, $H(A)$, is the entropy of the array's activation
pattern marginalized over sensing events. Index the $B$ sampled events by
$i=1,\dots,B$ and let $C$ be the (uniform) random variable indicating which
event generated a given sample of $A$. Conditional on a specific event,
receptor activations are independent Bernoulli variables with parameters
$A_{i,r}$ (Sec.~\ref{sec:env}), so $p(A\mid C=i)$ is a simple, closed-form
distribution. The unconditional distribution $p(A)=\frac1B\sum_i p(A\mid C=i)$
is a \emph{mixture} of these $B$ simple components. This is a different use
of the word from the ligand mixtures of Sec.~\ref{sec:env}, though each
component here is populated by one such physical mixture: it refers to
$p(A)$ itself being a mixture of distributions, one per sensing event, not
to the chemistry of any single sensing event.

By definition of mutual information,
\begin{equation}
H(A) = H(A\mid C) + I(A;C).
\label{eq:mixture_decomp}
\end{equation}
The first term is cheap: conditioning on $C=i$ leaves $R$ independent
Bernoulli variables, so
\begin{equation}
H(A\mid C) \equiv H_\text{cond} = \frac{1}{B}\sum_i \sum_r h_2(A_{i,r}),
\qquad
h_2(a) = -a\log_2 a - (1-a)\log_2(1-a).
\label{eq:hcond}
\end{equation}
The second term, $I(A;C)$, is exactly as hard to compute as $H(A)$ itself:
it asks how distinguishable the $B$ components are, which still requires
the full joint distribution. Kolchinsky and Tracey give certified upper and
lower bounds on it from pairwise affinities between components
alone~\cite{kolchinsky_estimating_2017}, avoiding the $2^R$ enumeration.

\begin{itemize}
\item \textbf{Lower bound.} The Bhattacharyya affinity between events $i$
and $j$,
\begin{equation}
\log\text{BC}(i,j) = \sum_r \log\!\Bigl(\sqrt{A_iA_j}+\sqrt{(1-A_i)(1-A_j)}\Bigr),
\label{eq:BC}
\end{equation}
gives
\begin{equation}
H_\text{KT}^\text{low} = H_\text{cond} - \frac{1}{B}\sum_i \log_2\!\Bigl(\frac{1}{B}\sum_j \text{BC}(i,j)\Bigr).
\label{eq:KTlow}
\end{equation}
\item \textbf{Upper bound.} Replacing the Bhattacharyya affinity by minus
the Kullback-Leibler (KL) divergence between the same product-Bernoulli
components,
\begin{equation}
-\text{KL}(i\|j) = \sum_r \Bigl[A_i\log A_j + (1-A_i)\log(1-A_j)\Bigr] - \sum_r\Bigl[A_i\log A_i+(1-A_i)\log(1-A_i)\Bigr],
\label{eq:KL}
\end{equation}
gives
\begin{equation}
H_\text{KT}^\text{up} = H_\text{cond} - \frac{1}{B}\sum_i \log_2\!\Bigl(\frac{1}{B}\sum_j e^{-\text{KL}(i\|j)}\Bigr).
\label{eq:KTup}
\end{equation}
\end{itemize}

Both bracket terms have the form $H_\text{cond}+(\text{estimate of }I(A;C))$.
Because $C$ is uniform over $B$ values, $I(A;C)\le H(C)=\log_2 B$ always,
and both pairwise-affinity sums above are constructed to land in
$[0,\log_2 B]$ for the same reason: the diagonal term, $\text{BC}(i,i)=1$
and $\text{KL}(i\|i)=0$, dominates when the off-diagonal affinities vanish,
capping the correction at $\log_2 B$. It is only this correction term that
is bounded by $\log_2 B$, not $H(A)$ itself: $H_\text{cond}$ grows with $R$
and is unrelated to $B$, so $H(A)$ can exceed $\log_2 B$ once $H_\text{cond}$
is large enough, even though $I(A;C)$ alone cannot.

$H_\text{KT}^\text{low} \le H(A) \le H_\text{KT}^\text{up}$ is a certified
bracket for any $B$. It tightens as the off-diagonal affinities
$\text{BC}(i,j)$, $i\ne j$, shrink towards $0$, that is, as different
sensing events drive distinguishable, non-overlapping subsets of receptors
rather than similar, overlapping ones. This is precisely what optimizing
$\{\mathbf v_u^\pm\}$ to decorrelate the receptors and sharpen their
thresholds (Sec.~\ref{sec:env}) drives the array toward. In Fig.~2 of the
main text, solid lines report $H_\text{KT}^\text{up}$ and dotted lines
$H_\text{KT}^\text{low}$.

\bibliographystyle{apsrev4-2}
\bibliography{smelling}

\end{document}
